\documentclass[../main.tex]{subfiles}

\begin{document}

\chapter{Related Work}\label{ch:related-work}

Existing research and implementations were examined to identify approaches relevant to the 
development of a monitoring system for an IdP operating at the application layer. 
The primary objective was to identify solutions that address the overarching goal of this 
thesis and to analyse their respective limitations and areas of applicability.

These related works provide valuable insights and a foundation for this thesis by identifying 
existing approaches, their areas of coverage, and gaps that remain relevant to the 
considered scenario. 
The most relevant studies and projects are presented and analysed in the following sections.

\section{Why Monitor from the Inside?}\label{sec:why-inside}

As discussed in the introduction~\ref{ch:introduction}, the proposed monitoring system operates 
directly within the application layer of the Shibboleth IdP. 
Its objective is to detect attacks based on events and behaviour observable from 
within the application.

The underlying concept is not new. 
The OWASP AppSensor project describes an approach in which intrusion detection and response 
capabilities are integrated directly into an application~\cite{OWASPAppSensor}. 
In contrast to network-based security mechanisms, application-layer monitoring can make use of 
information about the application's internal state and semantics.

A network firewall primarily controls network communication, while a network-based \acrfull{ids} 
has limited visibility into application-specific behaviour. 
A \acrfull{waf} can provide additional protection against generic application-layer attacks, such 
as SQL injection, but typically has limited knowledge of the internal logic and state of a specific 
application. 
Consequently, application-integrated detection mechanisms can complement network- and perimeter-based 
security controls by analysing application-specific events and behaviour~\cite{OWASPAppSensor}.

This application-level perspective also provides a basis for classifying the proposed system 
within an established security taxonomy.
In the following, the proposed system is classified using the taxonomy for self-protecting software 
systems introduced by Yuan and Malek~\cite{Taxonomy}. 
Two dimensions of this taxonomy are particularly relevant to the scope of this work: the 
\texttt{Architecture Level} and the \texttt{Self-Protection Level}.

The \texttt{Architecture Level} describes the layer at which a security mechanism operates and 
distinguishes between Network, Host, Software/Application, and Services. 
This dimension allows the observation position of different security mechanisms to be compared. 
Network-based \acrfull{ids}, such as Suricata and Snort, operate primarily at the network level, 
whereas \acrfull{edr} solutions operate at the host level. 
The proposed system operates at the Software / Application level, as it observes events directly 
within the Shibboleth IdP. 
The distinction between these systems therefore lies primarily in the layer at which 
security-relevant information becomes observable.

While the Architecture Level describes where the system operates, the \texttt{Self-Protection Level} 
describes the degree of self-protection according to the MAPE process of Monitor, Analyze, 
Plan, and Execute. 
The proposed system fully covers the \textit{Monitor & Detect} level by collecting and detecting 
security-relevant events within the IdP. 
The \textit{Analyze & Characterize} level is partially addressed through event correlation and 
severity classification. 
In contrast, \textit{Plan & Prevent} is outside the scope of this work. 
The system generates alerts based on detected events but does not autonomously modify the system 
state or take preventive countermeasures.

This limited level of self-protection is intentional and follows from the passive design 
principle of this work.
The proposed system follows the concept of a \textit{Looking Glass}, in which security-relevant 
events are made observable without interfering with the operation of the monitored system. 
Rieckers applies a similar principle to the passive observation of a 
federation infrastructure~\cite{EAPeduroam}, which is adopted by this work. 
Extending the system with automatic response mechanisms would contradict this passive 
observation principle. 
Furthermore, erroneous automated responses could negatively affect the availability or operation 
of the IdP and thereby potentially impact multiple SPs that rely on it.

The Looking Glass principle is also reflected in existing application-level 
monitoring approaches. 
One relevant architectural precedent is BlackWatch by Hall et al.~\cite{BlackWatch}.
BlackWatch applies a similar architecture to generic web applications. 
Detection points are placed directly within the application to identify security-relevant events. 
These events are collected by the separate BlackWatch application, where they are analysed using 
rule-based and reputation-based mechanisms. 
The results are then made available through a web-based reporting interface.

The separation between the monitored application and the analysis component allows the detection 
and analysis functionality to be decoupled from the application itself. 
A key challenge identified by Hall et al. is the integration and configuration of detection 
points within the application. 
Developers must identify suitable locations where suspicious behaviour can be detected and 
implement the corresponding detection logic~\cite{BlackWatch}.

\section{Why SIEM Functionality Is Not Enough}\label{sec:why-not-siem}

A \acrfull{siem} primarily operates on logs and security-relevant events. 
This is not a criticism, but a description of its basic operating principle. 
Two limitations follow from this principle. 
First, information that is not represented in an event or log is not available to the SIEM for 
analysis. 
Second, the SIEM generally relies on the information provided by its data sources and therefore 
cannot independently verify information that is incorrect or has been manipulated at the source.

These limitations do not diminish the value of SIEM systems, but they define the boundary within 
which they operate. 
The collection, storage, and correlation of log data used in this work therefore correspond to 
established SIEM functionality and do not constitute an original contribution. 
OWASP A09 explicitly recommends centralized logging and monitoring and identifies solutions such 
as the ELK stack as possible tools for addressing logging and monitoring 
vulnerabilities~\cite{OWASPA09}.

The question therefore shifts from how events are analysed to how their underlying sources can 
be verified. 
Wazuh is particularly relevant in this context, as it combines log analysis with file integrity 
monitoring and security configuration assessment, thereby addressing both event-based analysis 
and the verification of system state~\cite{WazuhDocumentation}. 
However, these mechanisms primarily address configuration and file state rather than the 
relationship between configuration and individual runtime events. 
Security configuration assessment evaluates system configurations against defined policies, 
while file integrity monitoring detects changes to monitored files. 
Neither mechanism directly establishes whether a particular runtime event is consistent with the 
configuration state that was active when the event occurred.

A further concern arises when considering the integrity of the log sources themselves. 
Graves demonstrates that log entries can be tampered with in a Shibboleth deployment and therefore 
recommends keeping Apache logs separate, since logs produced by Tomcat and the IdP itself can 
potentially be modified by an attacker with sufficient system access~\cite{ShibbolethGIAC}. 
OWASP A09 likewise identifies insufficient protection against log tampering as a security concern 
and refers to CWE-117, which describes the insertion, modification, deletion, or forgery of log 
entries to mislead subsequent analysis~\cite{OWASPA09,CWE117}.

This makes the correlation of independent event sources more than a technical convenience. 
It provides an additional means of identifying inconsistencies between records of the same event. 
A system that relies solely on log analysis can compare multiple log sources, but if all sources 
are stored as files on the same system, they remain exposed to an attacker with sufficient 
file-system access.

Cryptographic approaches provide a stronger solution to this problem. 
Schneier and Kelsey, as well as Snodgrass et al., describe mechanisms that make tampering with 
stored logs cryptographically detectable~\cite{TamperDetection,CryptoSupportSecureLogs}. 
However, these approaches require changes to the logging or storage architecture and may depend 
on external trusted components. 
Snodgrass et al. report an overhead of up to 16 percent for their approach~\cite{TamperDetection}.

The approach taken in this work is different. 
Instead of modifying the existing logging infrastructure, two independently generated records 
of the same event are compared. 
The difference is therefore not based on a new trust boundary, but on the temporal and 
architectural separation between the records. 
A log file may remain accessible to an attacker with file-system access for an extended period, 
whereas a report generated by the Custom Interceptor Flow exists only for the duration of the 
corresponding application processing. 
This does not make the interceptor report immutable: an attacker controlling the JVM process 
could also manipulate it. 
Its purpose is instead to provide a second, independently generated observation of the same 
event.

The remaining question is whether such application-level detection is already sufficiently 
covered by existing frameworks. 
Of the works considered, AppSensor is conceptually closest to this work. 
However, its detection model is based on different assumptions about the monitored 
application~\cite{OWASPAppSensor}.

AppSensor detection points primarily monitor the behaviour of a user within an application 
session. 
They can detect events such as tampered cookies, modified parameters, or injection attempts. 
The framework therefore asks whether a user's behaviour is suspicious. 
It does not provide a mechanism for comparing an observed runtime event with a configuration 
or security state declared by a third party.

This limitation follows from the architectural assumptions of the framework. 
AppSensor is designed for applications in which the application itself is the primary authority 
over the user's interaction. 
An IdP, in contrast, performs authentication and issues identity information on behalf of other 
entities. 
Consequently, security-relevant conditions may depend not only on the IdP's own runtime 
behaviour, but also on configuration and trust relationships involving its federation partners.

This creates a detection category that is not directly represented in the AppSensor model: 
comparing an observed runtime event against the configuration that should govern that event.
One AppSensor detection point, however, is directly applicable. 
IE6, \textit{Violation of Security Log Integrity}, addresses the integrity of security logs and 
cites NIST SP 800-92 and Snodgrass et al. as supporting references~\cite{NIST80092}. 
The discrepancy-based approach used in this work therefore does not introduce the general concept 
of detecting log integrity violations. 
Instead, it applies an established detection principle to the specific domain of federated 
identity management and extends it with application-level observations from the Shibboleth IdP.

\section{Derivation of the Rules}\label{sec:rules}

The derivation of the detection rules is based on three complementary sources: the FIM survey 
by Simpson and Groß, the OASIS SAML Security and Privacy Considerations, and the STRIDE threat 
classification scheme.
The FIM survey identifies 11 vulnerabilities and 14 attack classes based on a systematic 
literature review~\cite{FIMSurvey}. 
The vulnerabilities are grouped into cross-protocol issues where the same weakness is exploited 
by the same attack class, results in the same CIA consequence, and is documented for at least 
two different protocols. 
Six of the fourteen identified issues concern SAML, while one directly concerns Shibboleth. 
The survey therefore provides the primary basis for identifying attack scenarios that are 
relevant to federated identity management.

While the survey describes vulnerabilities that have been documented in the literature, the 
OASIS SAML Security and Privacy Considerations provide a normative perspective by specifying 
security requirements and recommended countermeasures for SAML 
deployments~\cite{SAMLSecurityConsideration}. 
This allows the detection rules to be derived not only from observed or documented attacks, 
but also from the security properties that a SAML deployment is expected to enforce.

A third perspective is provided by STRIDE, a threat classification scheme that categorizes 
threats according to their mode of operation~\cite{ThreatModeling}. 
It distinguishes between spoofing, tampering, repudiation, information disclosure, denial 
of service, and elevation of privilege. 
Graves applies this framework to the analysis of threats against 
Shibboleth~\cite{ShibbolethGIAC}. 
In this work, STRIDE is used to characterize the security impact of the attack scenarios 
identified by the other two sources.

The three sources therefore serve different purposes in the derivation process. 
The FIM survey identifies documented vulnerabilities and attack classes, OASIS specifies the 
security requirements that a SAML deployment should satisfy, and STRIDE provides a classification 
of the resulting security impact. 
Where the evidence from the survey and the requirements specified by OASIS overlap, the basis 
for a detection rule is particularly strong. 
Replay attacks provide an example of such an overlap, as they are documented in the literature 
and explicitly addressed by the SAML security considerations.

\section{What Is Structurally Impossible}\label{sec:structureally-impossible}

This section explains why the proposed monitoring system has some potential structurall 
difficulties. 
It demonstrates that the limitations of the approach have been thoroughly considered 
rather than merely listed.

Three Groups of Reasons:
- Other System Levels: 
    - Weak DNS and DNS poisoning affect name resolution. 
    - The same applies to DDoS and UI redressing
- Other party: 
    - Vulnerable SP and CSRF occur between the browser and the SP. 
    - The same applies to XSS and session swapping
- Only indirectly observable: 
    - Lack of Binding and MITM. 
    - The prerequisites are observable, not the event itself

Two special cases should be highlighted:
- No Trust Infrastructure and Bogus Merchant is the only Shibboleth-related entry in the 
survey and remains open. 
    - The malicious SP is registered normally and is formally legitimate.
- A malicious provider is structurally undetectable. 
    - Since the system runs within the IdP, the observer and the potential attacker are 
    one and the same. 
    - This limitation applies to every IdP-side system.

And now, the connection to motivation: 
- None of the thirteen rules reaches R5 presented in Graves.
- The Sunburst case from the beginning, specifically, would not be detectable 
with this system.

\end{document}